\subsection{Датчики дрона}

Все методы чтения датчиков устроены одинаково: они ничего не ждут, возвращают кортеж или число и возвращают \verb|None|, если данных нет (датчика нет на плате, нет связи, автопилот ещё не прислал значение). Поэтому результат всегда проверяют перед распаковкой. Взлетать для чтения датчиков не нужно.

\begin{center}
\small
\begin{tabularx}{\linewidth}{@{}>{\ttfamily\raggedright\arraybackslash}p{0.34\linewidth} X@{}}
\toprule
\textrm{\textbf{Метод}} & \textbf{Что возвращает} \\
\midrule
get\_accel() & ускорение $(a_x, a_y, a_z)$ по осям дрона \\
get\_gyro() & угловая скорость $(g_x, g_y, g_z)$, рад/с \\
get\_mag() & показания магнитометра $(m_x, m_y, m_z)$ \\
get\_orientation() & крен, тангаж, рыскание $(roll, pitch, yaw)$, градусы \\
get\_altitude() & высота, метры \\
get\_dist\_sensor\_data() & дальность оптического дальномера вниз, метры \\
get\_optical\_data() & оптический поток: $(flow_x, flow_y, дальность)$; поток --- целые числа датчика, дальность --- метры \\
get\_ranger\_data() & пять расстояний модуля Ranger в метрах: справа, слева, спереди, сзади, сверху/снизу \\
get\_motors\_rpm() & список из четырёх чисел --- обороты моторов \\
\bottomrule
\end{tabularx}
\end{center}

\begin{codefile}[python]{python-onboard/pioneer_sdk2/sdk2_08_sensors.py}
\end{codefile}

\paragraph{Что происходит.} Программа три раза подряд печатает все показания. Функция \verb|show| отличает «нет данных» (\verb|None|) от нуля. \verb|get_ranger_data()| устроен иначе: когда модуля Ranger нет, он возвращает не \verb|None|, а пять значений \verb|None| подряд (\verb|(None, None, None, None, None)|), поэтому распаковка проходит всегда, а проверять нужно первое значение. Если плата не поддерживает Ranger, библиотека дополнительно печатает «Модуль Ranger не поддерживается».

\paragraph{Какой датчик что даёт.}
\begin{itemize}
    \item \textbf{IMU} --- акселерометр, гироскоп, магнитометр. По акселерометру и гироскопу видно, как дрон движется и вращается; по магнитометру --- куда повёрнут нос относительно магнитного поля Земли. В описании пакета магнитометр указан «только для GPS»: на дроне с локальной навигацией он может вернуть \verb|None|. В симуляторе магнитометр моделирует поле Земли по курсу.
    \item \textbf{Ориентация} --- углы наклона и курса, посчитанные автопилотом. \verb|get_orientation()| даёт сразу три угла: крен, тангаж и рыскание; \verb|get_local_yaw_lps()| --- только курс.
    \item \textbf{Высота} --- \verb|get_altitude()| считает автопилот (в симуляторе --- по барометру); \verb|get_dist_sensor_data()| --- расстояние до поверхности под дроном по дальномеру; \verb|get_local_position_lps()[2]| --- высота в локальной системе координат. Это три разных способа, и значения у них не обязаны совпадать.
    \item \textbf{Оптический поток} --- нижняя камера и дальномер определяют, как «плывёт» пол под дроном; так дрон без Локуса и GPS удерживает позицию. Дальность \verb|get_optical_data()[2]| та же, что \verb|get_dist_sensor_data()|.
    \item \textbf{Ranger} --- модуль из пяти дальномеров; он измеряет расстояния в пяти направлениях (справа, слева, спереди, сзади, сверху или снизу). В симуляторе дальность ограничена 4~м.
    \item \textbf{Обороты моторов} --- четыре числа. Если не удалось прочитать хотя бы одно значение, в списке будут четыре \verb|None|.
\end{itemize}

\begin{infobox}[Единицы измерения]
SDK не указывает единицы у ускорения и магнитометра: это число из поля автопилота, поделённое на 1000. Гироскоп --- рад/с, ориентация --- градусы, высота и дальности --- метры. Чтобы узнать единицы ускорения самому, положите дрон ровно и посмотрите, чему равна ось, направленная вверх: это покажет, что за число вы измеряете.
\end{infobox}

\begin{taskbox}[Какие у ускорения единицы]
Оставьте неподвижный дрон на земле и напечатайте \verb|get_accel()| 10 раз. Определите, по какой оси значение заметно больше остальных, и по числу догадайтесь, в каких единицах ускорение.
\expected{все десять замеров близки друг к другу; одна из осей заметно отличается от остальных. В симуляторе значения могут быть модельными, а на настоящем дроне они покажут силу тяжести по этой оси.}
\end{taskbox}

\begin{taskbox}[Тревога по Ranger]
Если Ranger нет (первое значение \verb|None|), завершите программу с сообщением до взлёта. Иначе взлетите на 1~м и летите вперёд со скоростью 0.2~м/с короткими командами (\verb|interval=0.5|), пока \verb|get_ranger_data()| не покажет впереди меньше 0.5~м; тогда остановитесь (нулевая скорость) и сядьте.
\expected{перед препятствием дрон останавливается до столкновения и садится; без Ranger печатается сообщение и взлёта нет.}
\end{taskbox}

\subsection{Журнал полёта в CSV}

Показания датчиков удобно записывать в файл: по журналу видно, что происходило в полёте. Запись ведёт отдельный фоновый поток параллельно с миссией.

\begin{codefile}[python]{python-onboard/pioneer_sdk2/10_telemetry_logger.py}
\end{codefile}

\paragraph{Что происходит.} Класс \verb|TelemetryLogger| наследует \verb|threading.Thread|, логика записи --- в \verb|run()|. \verb|daemon=True| не даёт потоку удерживать программу после завершения основного кода. Остановка идёт через \verb|threading.Event|: \verb|wait(period)| работает как пауза, которую можно прервать, поэтому после \verb|stop()| поток завершается сразу. Имя \verb|_stop_event| выбрано не случайно: у \verb|threading.Thread| уже есть внутренний метод \verb|_stop|, и его переопределение ломает \verb|join()|. Запись \verb|value or [None] * 3| подставляет пустые значения, если метод вернул \verb|None|. Имя файла содержит дату и время, поэтому журналы разных полётов не перезаписывают друг друга.

Для анализа скопируйте CSV на компьютер и постройте графики (программа запускается на компьютере, нужен \texttt{matplotlib}):

\begin{codefile}[python]{python-onboard/pioneer_sdk2/10b_plot_log.py}
\end{codefile}

\begin{taskbox}[Журнал с событиями]
Добавьте в журнал колонку \texttt{event} с событиями автопилота (подписка из раздела про события) и колонку с текущей целевой точкой. Определите по журналу максимальную скорость перелёта (разность координат по времени).
\expected{в CSV видны строки с событиями \texttt{POINT\_REACHED} и \texttt{TAKEOFF\_COMPLETE}; скорость перелёта получается порядка десятых долей м/с.}
\end{taskbox}

\subsection{Навигация: LPS, GPS и оптика}

Автопилот определяет положение дрона одной из трёх систем навигации; какая используется --- решает параметр автопилота, а библиотека его читает:

\begin{codebox}[Система навигации]
\verb|get_nav_system(update=False)| --- возвращает \verb|NavSystem|: \verb|GPS| (спутники ГНСС), \verb|LPS| (локальная система «Локус» --- ультразвук и ИК-маяки) или \verb|OPT| (оптическая навигация по камере и дальномеру). Значение запоминается; \verb|update=True| читает заново.
\end{codebox}

Для каждой системы навигации в библиотеке свои методы чтения, и они работают только когда эта система установлена и выбрана:

\begin{center}
\small
\begin{tabularx}{\linewidth}{@{}>{\ttfamily\raggedright\arraybackslash}p{0.40\linewidth} X@{}}
\toprule
\textrm{\textbf{Метод}} & \textbf{Что возвращает} \\
\midrule
get\_nav\_status\_lps() & \verb|NavStatus| локальной навигации \\
get\_local\_position\_lps() & позиция $(x, y, z)$, метры \\
get\_local\_velocity\_lps() & скорость $(v_x, v_y, v_z)$, м/с \\
get\_local\_yaw\_lps() & курс, градусы, от $-180$ до $180$ \\
get\_nav\_status\_gps() & \verb|NavStatus| GPS \\
get\_global\_position\_gps() & $(широта, долгота, высота)$: градусы, градусы, метры \\
get\_global\_velocity\_gps() & $(v_{север}, v_{восток}, v_{вертикаль})$, м/с \\
get\_satellites\_count() & $(число спутников GPS, число спутников ГЛОНАСС)$ \\
\bottomrule
\end{tabularx}
\end{center}

Если модуль навигации не установлен, метод печатает «Модуль навигации не установлен» и возвращает \verb|None|; в \verb|get_global_*| и \verb|get_satellites_count| значение \verb|None| означает и «нет модуля», и «модуль не отвечает». Модуль ГНСС библиотека находит при создании объекта по списку поддерживаемых модулей платы.

\paragraph{Статус навигации.} \verb|NavStatus| показывает, можно ли верить позиции: \verb|NO_DATA| --- модуль ничего не прислал, \verb|CANNOT| --- позицию оценить невозможно, \verb|LOW| --- доверие малое, \verb|OK| --- оценка успешна. Взлетать надо при \verb|OK|: если \verb|arm()| не удаётся, первое, что стоит проверить --- статус навигации.

\begin{codefile}[python]{python-onboard/pioneer_sdk2/sdk2_09_safe_start.py}
\end{codefile}

\paragraph{Что происходит.} Программа проверяет систему навигации и её статус, и только потом запускает моторы. Состояние \verb|get_fly_state()| печатается на каждом этапе. Перелёт в точку отслеживает оба события: сначала \verb|point_deceleration()| (дрон начал тормозить), затем \verb|point_reached()|. После прибытия \verb|flight_time()| показывает, сколько секунд идёт полёт. Нет надёжной навигации --- \verb|SystemExit| до взлёта: в \verb|finally| тогда сработает только \verb|disarm()| и \verb|close_connection()|.

\begin{warnbox}[Скорость по GPS идёт в других осях]
\verb|get_global_velocity_gps()| возвращает скорость по сторонам света (север, восток, вертикаль), а не по осям локальной системы, как \verb|get_local_velocity_lps()|. Не подставляйте одно вместо другого. Знак вертикальной составляющей --- по названию поля автопилота, это скорость «вниз»; проверьте его на земле лёгким подъёмом, прежде чем использовать в расчётах.
\end{warnbox}

\begin{taskbox}[Готов ли дрон взлетать]
Напишите функцию \verb|nav_ok(pioneer)|: она печатает имя системы навигации и статус нужной системы (LPS --- \verb|get_nav_status_lps|, GPS --- \verb|get_nav_status_gps|; для OPT отдельного метода статуса нет, напечатайте, что статус неизвестен) и возвращает \verb|True|, только если имя статуса \verb|OK|. Используйте её перед \verb|arm()|.
\expected{в симуляторе печатается \texttt{LPS} и \texttt{OK}, функция возвращает \texttt{True}, дрон взлетает; при статусе, отличном от \texttt{OK}, функция возвращает \texttt{False} и взлёта нет.}
\end{taskbox}

\begin{taskbox}[Спутники и позиция]
На дроне с модулем GPS (на улице) напечатайте раз в секунду 10 раз: число спутников, статус и широту с долготой с шестью знаками. Дождитесь, когда статус станет \verb|OK|. В симуляторе GPS нет --- убедитесь, что методы возвращают \verb|None| и программа не падает.
\expected{на улице число спутников растёт, статус доходит до \texttt{OK}, координаты перестают «плавать»; в симуляторе все строки --- «нет данных», без ошибок и исключений.}
\end{taskbox}

\subsection{Параметры автопилота}

У автопилота есть настройки --- параметры: имя и число. Через SDK их можно читать и менять:

\begin{codebox}[get\_param и set\_param]
\verb|get_param(name, update=False)| --- значение параметра \verb|name| (число) или \verb|None|, если такого параметра нет или не удалось прочитать. При первом обращении библиотека загружает список параметров автопилота и печатает «-> Загрузка словаря автопилота ...»; \verb|update=True| загружает список заново.\\
\verb|set_param(name, value)| --- записать значение (оно хранится как вещественное число). Возвращает \verb|True|, если параметр с таким именем есть и команда отправлена, иначе \verb|False|.
\end{codebox}

Именно так \verb|get_nav_system()| узнаёт навигацию: читает параметр \texttt{Flight\_com\_navSystem}. Высоту взлёта задаёт параметр \texttt{Copter\_com\_takeoffAlt}. Описание параметров смотрите в документации автопилота вашей платы. Чтобы получить свежее значение из автопилота, а не запомненное при включении, вызывайте \verb|get_param(name, update=True)|.

\begin{codebox}[Безопасная смена параметра]
\small
\begin{verbatim}
name = "Copter_flyWithoutRc"
old = pioneer.get_param(name)
print("было:", old)
if old is not None:
    new = 1.0          # ваше значение; здесь можно взять и old
    pioneer.set_param(name, new)
    print("стало:", pioneer.get_param(name))
    # ... работа ...
    pioneer.set_param(name, old)   # вернуть как было
\end{verbatim}
\end{codebox}

\begin{warnbox}[Параметры меняют поведение дрона]
Опечатка в имени параметра безвредна (\verb|set_param| вернёт \verb|False|), а вот верное имя с неверным значением изменит поведение дрона в полёте. Перед записью читайте старое значение, запоминайте его и возвращайте в конце программы. По описанию Pioneer Code параметры автопилота \emph{сохраняются после перезагрузки}, поэтому ошибочное значение (например, источника сигнала RC) может сделать дрон неуправляемым с пульта. В симуляторе параметров нет: \verb|get_param| возвращает \verb|None|.
\end{warnbox}

\paragraph{Контекстный менеджер для параметров.} Возврат старого значения удобно спрятать в конструкцию \verb|with|:

\begin{codefile}[python]{python-onboard/pioneer_sdk2/11_params_context.py}
\end{codefile}

Код до \verb|yield| выполняется при входе в блок \verb|with|, код в \verb|finally| --- при выходе, \emph{даже если внутри блока была ошибка}. Значения читаются с \verb|update=True|, результат \verb|set_param| проверяется. Программа временно ставит высоту взлёта 0.7~м и возвращает прежнюю. В симуляторе параметров нет (\verb|get_param| возвращает \verb|None|), и пример там остановится на проверке записи: запускайте его на дроне.

\begin{taskbox}[Резервная копия параметров]
Напишите функцию \verb|with_param(pioneer, name, value, action)|, которая читает старое значение параметра, записывает новое, выполняет \verb|action()| и в \verb|finally| возвращает старое значение. Проверьте её на неверном имени параметра: если параметра нет, действие выполнять не нужно. Для проверки пишите в параметр то же значение, что прочитали.
\expected{для верного имени параметр после программы равен исходному значению, даже если \texttt{action()} упал; для неверного имени \texttt{set\_param} печатает «Установка параметра ... не удалась»; ваша функция в этом случае не выполняет действие и сообщает об ошибке.}
\end{taskbox}
